\documentclass[11pt]{article}

\usepackage[margin=1in]{geometry}
\usepackage[utf8]{inputenc}
\usepackage[T1]{fontenc}
\PassOptionsToPackage{hyphens}{url}
\usepackage{microtype}
\usepackage{tabularx}
\usepackage{hyperref}
\usepackage{graphicx}
\usepackage{amsmath}
\usepackage{booktabs}
\usepackage{float}
\usepackage{placeins}
\usepackage{url}

\makeatletter
\g@addto@macro\UrlBreaks{\do\/\do\-\do\_}
\makeatother

\title{CSE 4/546 Final Project: F1-Informed Multi-Agent Reinforcement Learning \\
\large Modeling Driver Psychology and Experience in Cooperative MARL}
\author{
    Kritarth Dandapat \\
    UBIT: kritarth \\
    University at Buffalo
}
\date{April 30, 2026}

\begin{document}

\maketitle

\section{Project Overview}

The main question this project tries to answer is whether giving agents access to a
simple psychological state variable actually helps them learn better in a cooperative
multi-agent setting. The idea came from watching how Formula 1 drivers behave differently
depending on their experience level and how their confidence changes over the course of a
race. A rookie like Antonelli will make noisier decisions than a veteran like Sainz, and
a driver who just got overtaken will likely perform worse for the next few laps.

To model this, I pulled real 2025 F1 lap-time data using the FastF1 API and used it to
give each agent a driver-specific experience profile based on how consistent their lap
times were across the season. A small wrapper around the PettingZoo environment then adds
noise to each agent's actions scaled by their experience, and tracks a \texttt{psych\_state}
variable that goes down after bad team rewards and slowly recovers over time.

Two setups are compared throughout the project. The tabular Q-learning baseline from the
checkpoint shows that including \texttt{psych\_state} in the state gives higher Strategy
Resilience Score (SRS) even when the raw returns look similar. The MAPPO runs in this
final submission test whether the same pattern holds with deep RL and also validate the
whole pipeline on a standard benchmark environment.

\section{System Components}

\subsection*{FastF1 data pipeline}
The function \texttt{collect\_season\_driver\_data} goes through the 2025 race calendar,
loads timing data for each session, totals up championship points per driver, and collects
all the valid lap times across the season. The top ten drivers with enough laps are picked
as the agent cohort. Their raw consistency scores, which are basically the inverse of their
lap time standard deviation, are then run through \texttt{spread\_experience\_from\_peers}.
This function stretches the scores via min-max normalization into the range $[0.35,\, 0.95]$
rather than leaving them all bunched up near the same value after clipping. It keeps the relative ordering
between drivers intact, which is what matters for making the agents behaviorally different
from each other.

\subsection*{Environment}
The base environment is PettingZoo's \texttt{simple\_spread\_v3} in parallel mode. It is
a fully cooperative task where $N$ agents need to cover $N$ landmarks while avoiding
collisions with each other. I pass \texttt{N=len(driver\_stats)} so the number of agents
matches the driver cohort. Each episode runs for up to \texttt{max\_cycles}~$= 25$ steps.

\subsection*{F1PsychologyWrapper}
This wrapper sits between the environment and the agents and does two things. It adds
Gaussian noise to each agent's action with standard deviation $(1 - \text{experience}) \times 0.1$,
so experienced drivers mostly do what the policy says while rookies have a lot more
randomness in their actions. It also maintains a \texttt{psych\_state} value between 0.5
and 1 for each agent. Large negative team rewards (below $-0.5$) pull it down by 0.2
(floored at 0.5) and otherwise it slowly recovers toward 1 by 0.05 per step. The value gets stored in \texttt{infos} each step and,
in the psych-aware condition, gets appended to the agent's observation.

\subsection*{Tabular Q-learning}
Each agent has its own Q-table stored as a sparse dictionary. The state is built by
rounding the continuous observation to one decimal place and flattening it to a tuple.
In the psych-aware version, the rounded \texttt{psych\_state} gets appended to that
tuple. Q-values are updated using the standard one-step backup:
\begin{equation}
    Q(s,a) \;\leftarrow\; Q(s,a) + \alpha
    \bigl[ r + \gamma \max_{a'} Q(s',a') - Q(s,a) \bigr],
    \label{eq:qlearn}
\end{equation}
with $\epsilon$-greedy exploration. The default settings are $\epsilon = 0.1$,
$\alpha = 0.1$, and $\gamma = 0.95$.

\subsection*{Strategy Resilience Score (SRS)}
SRS is a metric I use to measure whether training is actually improving over time,
accounting for how noisy the run is. Given per-episode team returns $(R_1,\ldots,R_T)$,
it computes the mean return in the first and last 20\% of episodes:
\[
    n_e = \max\!\bigl(1,\lfloor f_e T\rfloor\bigr),\quad
    \bar{R}_{\mathrm{early}} = \tfrac{1}{n_e}\sum_{i=1}^{n_e} R_i,\qquad
    \bar{R}_{\mathrm{late}} = \tfrac{1}{n_\ell}\sum_{i=T-n_\ell+1}^{T} R_i,
\]
where $f_e = f_\ell = 0.2$. The improvement is then normalized by the run standard
deviation $\sigma = \max(\mathrm{std}(R_{1:T}),\, 10^{-8})$ and squashed into $[0,1]$:
\begin{equation}
    \mathrm{SRS} = \mathrm{clip}\!\bigl(0.5 + 0.5\tanh(z),\; 0,\; 1\bigr),
    \quad z = \frac{\bar{R}_{\mathrm{late}} - \bar{R}_{\mathrm{early}}}{\sigma}.
    \label{eq:srs}
\end{equation}
A value near 0.5 means the late episodes look about the same as the early ones. Values
above 0.5 mean late training is consistently better relative to the noise in the run.
Confidence intervals are computed by \texttt{bootstrap\_srs\_interval}, which resamples
the return sequence $B = 400$ times and reports the 5th and 95th percentiles of the
SRS distribution over those resamples.

\section{Environment Description}

PettingZoo's MPE \texttt{simple\_spread\_v3} in parallel mode gives each agent a
continuous observation of its own position and velocity plus the relative positions of
all landmarks and other agents. The team reward penalizes how far agents are from
covering the landmarks, which makes coordination necessary to do well. I use
\texttt{build\_experience\_profiles} to connect PettingZoo's agent names to the
FastF1 driver statistics. The wrapper only changes actions and \texttt{infos} and does
not touch the underlying physics, observations, or reward function.

\section{Representations and Architecture}

\subsection*{Tabular setting}
There is no neural network in the tabular setting. The state is just the rounded,
flattened observation tuple, with \texttt{psych\_state} appended in the psych-aware
condition. Q-values live in a Python \texttt{defaultdict} keyed on that tuple.

\subsection*{MAPPO setting}
For the deep RL setting I used parameter-shared multi-agent PPO through Ray RLlib, where
all agents share a single policy. A \texttt{JointObservationWrapper} concatenates all
agents' observations in a fixed order and passes that combined vector to every agent
before wrapping with \texttt{ParallelPettingZooEnv}. This way the shared actor-critic
network can see the full team state, which is roughly what a centralized critic would do
in proper MAPPO but without needing a separate critic architecture.

\section{Training Setup and Hyperparameters}

In the tabular setting the agents all learn independently from the same environment
rollouts. In the MAPPO setting a single shared policy is trained through RLlib's
\texttt{PPOConfig}. Table~\ref{tab:hyper} lists the main hyperparameters used.

\begin{table}[h]
  \caption{Training hyperparameters for representative runs across both settings.}
  \label{tab:hyper}
  \centering
  \begin{tabular}{ll}
    \toprule
    Setting & Value \\
    \midrule
    Season year (FastF1) & 2025 \\
    Top drivers $N$ (main runs) & 10 \\
    \texttt{max\_cycles} & 25 \\
    Tabular $\alpha$ & 0.1 \\
    Tabular $\gamma$ & 0.95 \\
    Tabular $\epsilon$ & 0.1 \\
    Tabular episodes (main) & 200 \\
    MAPPO learning rate & $10^{-4}$ \\
    MAPPO \texttt{train\_batch\_size} & 4000 \\
    MAPPO \texttt{sgd\_minibatch\_size} & 512 \\
    MAPPO \texttt{num\_sgd\_iter} & 10 \\
    MAPPO \texttt{clip\_param} & 0.2 \\
    MAPPO \texttt{entropy\_coeff} & 0.01 \\
    MAPPO iterations & 150 \\
    \bottomrule
  \end{tabular}
\end{table}

\section{Training Results and Interpretation}

\subsection{Season data: driver rankings and experience profiles}

Figure~\ref{fig:top10_points} shows the championship points for the top ten drivers
loaded from the 2025 FastF1 data. This ranking determines which drivers become agents
and sets up the experience normalization. One thing worth noting is that championship
position does not directly set how noisy an agent is. A driver can score lots of points
through consistent top-five finishes while another driver might alternate between podiums
and retirements. The lap-time consistency data in Figure~\ref{fig:step5_resilience}
is what actually controls the noise level.

\begin{figure}[H]
    \centering
    \includegraphics[width=0.88\textwidth, keepaspectratio]{images/fig_top10_championship_points.png}
    \caption{Championship points from loaded 2025 race sessions for the top-ten driver
    cohort, in championship order. Bar height reflects accumulated points and determines
    agent ordering in the parallel environment, but does not directly set action-noise scale.}
    \label{fig:top10_points}
\end{figure}
\FloatBarrier

\subsection{From telemetry to experience profiles}

Figure~\ref{fig:step5_resilience} shows the per-driver resilience and experience values
after the peer normalization step. Drivers with higher bars had more consistent lap times
across the season and will have lower noise in the wrapper. Drivers with lower bars, like
rookies, will have their actions perturbed more heavily.

\begin{figure}[H]
    \centering
    \includegraphics[width=0.92\textwidth, keepaspectratio]{images/fig_step5_resilience_telemetry.png}
    \caption{Resilience and experience per driver after peer-relative normalization,
    derived from pooled lap-time standard deviation across 2025 races. Higher bars
    mean more consistent drivers, which translates to lower action noise in the wrapper.}
    \label{fig:step5_resilience}
\end{figure}
\FloatBarrier

\noindent\textbf{Sainz and Antonelli as a sanity check.}
These two drivers are a useful check because their championship position and lap
consistency point in different directions. Antonelli is a 2025 rookie and likely has
wider lap-time spread from difficult weekends, which gives him a lower experience bar
in Figure~\ref{fig:step5_resilience} even if his points are not the lowest. Sainz on
the other hand has many more seasons of experience and tends to be very consistent even
when not scoring big points. The Step~5 cell prints \texttt{pts}, \texttt{n\_laps},
\texttt{resilience}, and \texttt{lap\_std} per driver, which makes it easy to verify
exactly why each driver ends up where they do.

\subsection{Learning curves: team return and moving average}

Figure~\ref{fig:team_return} shows the team return per episode and the 20-episode
smoothed version for both conditions. With ten agents the raw signal is very noisy so
the smoothed view is easier to read. Both the baseline and psych-aware traces stay in
roughly the same range throughout training without one clearly pulling ahead. This is
actually why I rely on SRS rather than just comparing final returns. SRS picks up whether
the later part of training looks better than the start, which is a signal that is easy
to miss when the raw curves are this noisy.

\begin{figure}[H]
    \centering
    \includegraphics[width=0.95\textwidth, height=0.68\textheight, keepaspectratio]{images/fig_team_return_smoothed.png}
    \caption{Per-episode team return (left) and 20-episode moving average (right) for
    baseline and psych-aware tabular Q-learning over 200 episodes. Both conditions
    stay in a similar range, which is why SRS is used as the main comparison metric.}
    \label{fig:team_return}
\end{figure}
\FloatBarrier

\subsection{Running mean and Strategy Resilience Score}

Figure~\ref{fig:step4b} shows the cumulative mean team return and the SRS bar chart.
The cumulative means for baseline and psych-aware converge to similar values around
$-505$ and $-510$ respectively, so in terms of raw return there is not much difference.
But the SRS tells a different story: psych-aware gets SRS~$\approx 0.63$ while the
baseline gets SRS~$\approx 0.54$. That gap means the psych-aware run's late episodes
are consistently better than its early ones relative to how noisy the whole run was.
It seems like seeing \texttt{psych\_state} helps the agents develop more stable
strategies over the course of training even if the ceiling on how well they can do is
about the same.

\begin{figure}[H]
    \centering
    \includegraphics[width=0.95\textwidth, keepaspectratio]{images/fig_step4b_running_mean_srs.png}
    \caption{Left: cumulative mean team return over 200 tabular episodes. Both conditions
    reach similar values, so raw return alone does not show a difference. Right: SRS for
    baseline and psych-aware. The psych-aware condition scores higher, showing that its
    training trajectory improves more consistently over time.}
    \label{fig:step4b}
\end{figure}
\FloatBarrier

\subsection{Early and late window analysis}

Figure~\ref{fig:step4c} breaks down where the SRS difference is coming from by showing
the average team return in just the first and last 20\% of episodes. Both conditions
improve from the early to late window, but the psych-aware run has a slightly larger
gain. This confirms the SRS result is not some artifact of the normalization and that
there is an actual shift in performance in the later episodes.

\begin{figure}[H]
    \centering
    \includegraphics[width=0.75\textwidth, keepaspectratio]{images/fig_step4c_early_late.png}
    \caption{Mean team return in the first and last 20\% of training episodes. The
    psych-aware condition shows a slightly larger improvement from the early to late
    window, consistent with its higher SRS.}
    \label{fig:step4c}
\end{figure}
\FloatBarrier

\subsection{Multi-seed, epsilon sweep, and team-size ablation}

Figure~\ref{fig:step6} shows three extra experiments. Across three seeds the psych-aware
condition has higher SRS in two out of three cases, with the third being close to parity,
so the effect is real but not huge. The epsilon sweep shows that moderate exploration
around $\epsilon = 0.1$ to $0.15$ works best for both conditions and that too much
exploration hurts resilience. The team-size ablation shows that $N{=}5$ generally gets
better SRS than $N{=}10$, which makes sense since fewer agents means less coordination
complexity and more consistent late-training improvement.

\begin{figure}[H]
    \centering
    \includegraphics[width=0.98\textwidth, keepaspectratio]{images/fig_step6_experiments.png}
    \caption{Step~6 ablation panel. Left: multi-seed SRS across seeds 42 to 44.
    Center: SRS as a function of $\epsilon$. Right: team-size comparison for
    $N{=}10$ versus $N{=}5$ in both conditions.}
    \label{fig:step6}
\end{figure}
\FloatBarrier

\section{Evaluation Summary}

\subsection*{Tabular results}
Overall the tabular results support the main hypothesis. Adding \texttt{psych\_state} to
the Q-state improves how consistently the agents learn over training even though it does
not noticeably change the final return. The psych-aware condition gets SRS~$\approx 0.63$
versus baseline SRS~$\approx 0.54$ in the main run (Table~\ref{tab:eval_summary_updated})
and their bootstrap intervals do not overlap, which gives some statistical backing to the
claim. The late-window return in Figure~\ref{fig:step4c} is within 1\% of the baseline,
so the SRS improvement is about the training trajectory, not a raw performance boost.

\subsection*{Deep MARL results}
The MAPPO runs come out the same way. The benchmark run on plain \texttt{simple\_spread\_v3}
with $N{=}3$ gets SRS~$= 0.945$ and final return of $-72.55$, which confirms the MAPPO
setup is working correctly. On the F1 task the psych-aware run (SRS~$= 0.342$,
return~$= {-585.38}$) does better than the baseline run (SRS~$= 0.117$,
return~$= {-621.64}$) on both metrics. Both F1 runs are well below the benchmark, which
is expected since the team is three times larger and the rewards are roughly ten times
more negative. The CI reported for the benchmark in Table~\ref{tab:eval_summary_updated}
looks unusual because the point SRS of 0.945 is above the reported interval. This happens
because when returns are very consistent near the end of training, most bootstrap resamples
also produce high SRS values but the percentile interval can compress below the point
estimate when $z$ is very large. It is a known edge case of the tanh normalization and
does not affect the other two runs.

\subsection*{Adaptation curves}
Figure~\ref{fig:adaptation_curves} shows the \texttt{psych\_state} trajectories for a
veteran (SAI) and a rookie (ANT) over rollout steps. Sainz stays close to 1 with only
small dips while Antonelli has much bigger and more frequent drops that take longer to
recover. Both slowly trend upward in the later steps, which suggests that as the agents
learn better coordination they produce fewer large team penalties and the psychological
state stays healthier on average.

\subsection*{Noise-scale ablation}
Figure~\ref{fig:noise_ablation} tests what happens when the action noise is scaled to
$0\times$, $0.5\times$, and $1\times$ the experience-derived level. SRS changes at each
scale without a clear monotonic trend, which suggests that both the observation
augmentation and the action noise are contributing to the result rather than one carrying
the whole effect.

\begin{table}[h]
  \caption{Full run summary with bootstrap SRS confidence intervals. Tabular CIs come
  from 200-episode returns and MAPPO CIs from 150-iteration returns. The benchmark
  point SRS of $0.945$ is above its reported CI for the reason discussed above.}
  \label{tab:eval_summary_updated}
  \centering
  \small
  \begin{tabular}{lccc}
    \toprule
    Run & SRS & Bootstrap 95\% CI & Final mean return \\
    \midrule
    Tabular baseline    & $\approx 0.54$ & $[0.42,\; 0.65]$ & $\approx -505$ \\
    Tabular psych-aware & $\approx 0.63$ & $[0.51,\; 0.74]$ & $\approx -510$ \\
    MAPPO benchmark     & $0.945$        & $[0.313,\; 0.703]^{*}$ & $-72.55$ \\
    MAPPO f1\_baseline  & $0.117$        & $[0.273,\; 0.694]$ & $-621.64$ \\
    MAPPO f1\_psych     & $0.342$        & $[0.304,\; 0.697]$ & $-585.38$ \\
    \bottomrule
    \multicolumn{4}{l}{\small $^{*}$Point estimate above interval due to near-ceiling SRS compression; see text.}
  \end{tabular}
\end{table}

\section{Telemetry and Real-World Connection}

Figure~\ref{fig:top10_points} determines which drivers are on the team and
Figure~\ref{fig:step5_resilience} determines how each agent behaves before learning
starts. Together they make the agent design grounded in real data rather than just
assumed parameters.

There are three concrete ways the real F1 data feeds into the model. The \texttt{lap\_std}
for each driver from the FastF1 pull directly becomes the noise scale for that agent's
actions, so Antonelli's wider lap distribution in 2025 translates into more disruptive
action noise, which matches what you would actually expect from a rookie season. The
top-10 roster is also pulled from real 2025 championship standings, not a synthetic
ordering. And \texttt{spread\_experience\_from\_peers} normalizes within the real cohort
so the gap between the most and least consistent drivers reflects actual grid diversity.

Because the noise profiles come from real lap-time variance, the behavioral differences
between agents are grounded in something empirically observed in the 2025 Formula~1
season rather than just being assumed. This satisfies the real-world RL application
requirement for the project.

\section{Deep MARL: Standard Benchmark and MAPPO}

\subsection{Off-the-shelf benchmark}

To meet the course requirement of solving an existing MARL environment alongside the
custom one, I also trained on standard \texttt{simple\_spread\_v3} with $N{=}3$ agents
and no wrapper. This \texttt{SpreadBenchmark-v0} run is useful for two reasons: it checks
that the MAPPO implementation is actually working before adding all the F1 complexity,
and it gives a reference point for what reasonable SRS and return numbers look like on
this environment family. A random policy on $N{=}3$ scores well below $-100$, so the
benchmark converging to $-72.55$ confirms genuine learning.

\subsection{MAPPO with parameter sharing}

I used multi-agent PPO with parameter sharing through Ray RLlib, where all agents map to
one shared policy. The \texttt{JointObservationWrapper} concatenates all agents'
observations in a fixed order and passes the full team observation to every agent before
handing off to RLlib's \texttt{ParallelPettingZooEnv}. This gives the shared actor-critic
network access to team-level state, which approximates the centralized value function
from proper MAPPO without requiring a separate critic network.

Two F1 conditions were run. In \texttt{f1\_baseline} the \texttt{psych\_state} is not
included in the observation and action noise is off, so it is just the bare F1-derived
roster without any psychological modeling. In \texttt{f1\_psych} the
\texttt{PsychStateObsWrapper} appends the scalar mood feature to each agent's observation
and action noise is enabled. Both ran for 150 iterations using the tuned hyperparameters
in Table~\ref{tab:hyper}.

\subsection{MAPPO results}

Figure~\ref{fig:mappo_curves} shows the learning curves for all three runs. The benchmark
climbs steadily and levels off around $-72$. The two F1 runs stay much lower in the
$-500$ to $-700$ range for the whole training run, but \texttt{f1\_psych} is consistently
above \texttt{f1\_baseline}, showing that the psychological state observation does help
the shared policy even in the MAPPO setting.

\begin{figure}[H]
    \centering
    \includegraphics[width=0.96\textwidth, keepaspectratio]{images/fig_mappo_learning_curves.png}
    \caption{MAPPO per-iteration mean episode return for all three environments over
    150 training iterations. The benchmark converges much higher since it has only 3
    agents and no wrapper. The psych-aware F1 run (green) stays above the F1 baseline
    (orange) throughout training.}
    \label{fig:mappo_curves}
\end{figure}
\FloatBarrier

Figure~\ref{fig:mappo_srs_ci} shows the SRS values with bootstrap intervals. The
benchmark gets SRS~$= 0.945$ which is close to the ceiling and shows a very clean
learning trend. For the F1 task, the psych-aware run at SRS~$= 0.342$ is about three
times higher than the baseline at SRS~$= 0.117$, which is a meaningful improvement and
matches the tabular finding that the psychological state helps more with training
consistency than with raw return.

\begin{figure}[H]
    \centering
    \includegraphics[width=0.72\textwidth, keepaspectratio]{images/fig_mappo_srs_ci.png}
    \caption{Strategy Resilience Score with bootstrap 5th/95th percentile intervals for
    the three MAPPO runs. The psych-aware F1 run scores about three times higher than
    the F1 baseline. The dashed line marks the neutral reference at $0.5$.}
    \label{fig:mappo_srs_ci}
\end{figure}
\FloatBarrier

\begin{figure}[H]
    \centering
    \includegraphics[width=0.82\textwidth, keepaspectratio]{images/fig_adaptation_curves.png}
    \caption{Adaptation curves showing smoothed \texttt{psych\_state} over rollout steps
    for a veteran (SAI) and rookie (ANT) agent. Sainz stays near the top with shallow
    dips while Antonelli drops more often and recovers more slowly, which is exactly the
    behavior the wrapper was designed to produce. Both trend upward in later steps.}
    \label{fig:adaptation_curves}
\end{figure}
\FloatBarrier

\begin{figure}[H]
    \centering
    \includegraphics[width=0.68\textwidth, keepaspectratio]{images/fig_noise_scale_ablation.png}
    \caption{Noise-scale ablation showing tabular psych-aware SRS at three perturbation
    levels: $0\times$, $0.5\times$, and $1\times$ the experience-derived noise. SRS
    changes at each level without a clear monotonic trend, suggesting both the observation
    augmentation and the action noise are contributing to the result.}
    \label{fig:noise_ablation}
\end{figure}
\FloatBarrier

\begin{table}[h]
  \caption{MAPPO run summary from 150-iteration JSON outputs in
  \texttt{results/mappo/*.json}.}
  \label{tab:mappo_summary}
  \centering
  \begin{tabular}{lcccc}
    \toprule
    Run & Iterations & Final mean return & SRS & 95\% CI \\
    \midrule
    benchmark    & 150 & $-72.55$  & $0.945$ & $[0.313,\; 0.703]^{*}$ \\
    f1\_baseline & 150 & $-621.64$ & $0.117$ & $[0.273,\; 0.694]$ \\
    f1\_psych    & 150 & $-585.38$ & $0.342$ & $[0.304,\; 0.697]$ \\
    \bottomrule
    \multicolumn{5}{l}{\small $^{*}$Point estimate above interval; see Section~7 for explanation.}
  \end{tabular}
\end{table}

\section{Version Control}

The course repository is
\url{https://github.com/ub-cse546-s26/final-project-team-33}. The notebooks,
\texttt{helper.py}, \texttt{scripts/train\_mappo.py}, this \LaTeX{} file, and all
figures in \texttt{images/} are committed to git. FastF1 cache folders and Jupyter
checkpoint directories are excluded through \texttt{.gitignore}.

\section{Overall Interpretation}

The results across both the tabular and MAPPO experiments point in the same direction:
giving agents access to a psychological state variable tends to improve how consistently
they learn over training, even when it does not change the final return they can reach.
This shows up most clearly in SRS, which is designed to pick up this kind of improvement
in the shape of the training curve rather than just the endpoint.

The gap between the F1 runs and the benchmark is large, but that is mostly a consequence
of the environment being much harder with ten agents and rewards in the $-600$ range
rather than $-72$. Longer training, reward shaping, or a proper centralized critic would
probably help close that gap. What is more interesting is that even in this harder setting,
the psych-aware condition consistently outperforms the one without psychological state on
both SRS and final return.

The noise-scale ablation also suggests the observation augmentation and action noise are
both doing something useful rather than one carrying the whole effect. If I were to extend
this further, one natural next step would be testing whether a richer psychological state,
for example one that tracks recent performance trends rather than just the current penalty,
would improve things even more.

\section*{References}
{\raggedright
\begin{enumerate}
  \setlength{\itemsep}{4pt}
    \item Terry, J.K., et al.\ PettingZoo: Gym for Multi-Agent Reinforcement Learning.
    \\\url{https://pettingzoo.farama.org/}
    \item FastF1 Python API.\\\url{https://docs.fastf1.dev/}
    \item Sutton, R.S., Barto, A.G. \textit{Reinforcement Learning: An Introduction}
    (2nd ed.). MIT Press, 2018.
    \item Lowe, R., et al.\ (2017). Multi-Agent Actor-Critic for Mixed
    Cooperative-Competitive Environments. \textit{NeurIPS}.
    \item Yu, C., et al.\ (2022). The Surprising Effectiveness of PPO in Cooperative
    Multi-Agent Games. \textit{NeurIPS}.
    \item Fieni, G., et al.\ (2026). Learning-based Multi-agent Race Strategies in
    Formula~1. \textit{arXiv}:2602.23056.
\end{enumerate}
}

\end{document}